% SPDX-License-Identifier: Apache-2.0
% 01_intro.tex: Introduction
\section{Introduction}
\label{sec:intro}

\IEEEPARstart{E}{very} design document needs pictures. A reader takes in a
block diagram faster than three paragraphs that say the same thing. The
trouble is that a picture is usually drawn once, in a drawing program, and
saved as an image. The design keeps changing after that, and the picture
does not. A year later, the figure shows a register that was removed and
leaves out the one that replaced it.

Blockpic takes a different approach. A diagram is a short text file, kept in
the repository beside the HDL it describes. The build draws the picture from
that file every time. Blockpic can also draw a picture from the netlist that
synthesis writes, so there is no file to forget at all. When the HDL
changes, the next build redraws the figure from the new netlist.

This document is its own example. Both of its figures were drawn by the
Bazel build that typeset it \cite{bazel}.

\subsection{The running example}

The example design is a 4-bit counter. It adds one to a register on every
rising edge of \code{clk}, clears the register while \code{reset} is high,
and shows the register on four LEDs. Listing~\ref{lst:counter} gives the
whole design.

\lstinputlisting[style=vlog,firstline=4,caption={The running example,
\code{counter.v}.},label={lst:counter},float=t]{counter.v}

The counter is small on purpose. Its schematic after synthesis still fits a
column, and every cell in that schematic can be traced back to a line of the
listing.

\subsection{How to read this document}

Section~\ref{sec:by-hand} draws the counter by hand, in the blockpic
language. Section~\ref{sec:from-netlist} draws the same counter from the
netlist that synthesis writes for an AMD 7-series part. The build that ties
the two to this PDF comes next (\S\ref{sec:build}). The document ends with
what blockpic does not yet do well (\S\ref{sec:limits}).
